\section{Resultados e Discussão}
\label{sec:resultados}

Esta seção apresenta e discute os resultados obtidos a partir da execução sistemática dos experimentos descritos na Seção~\ref{sec:metodologia} e implementados conforme detalhado na Seção~\ref{sec:desenvolvimento}. Os dados aqui analisados refletem o comportamento das quatro estratégias de controle populacional avaliadas — \textbf{substituição geracional}, \textbf{elitismo}, \textbf{população fixa simples} e \textbf{controle por idade (GAVaPS)} — ao longo de múltiplas execuções independentes do algoritmo genético aplicado ao problema das \textit{n}-rainhas.

A análise foca em cinco aspectos principais:

\begin{enumerate}
    \item taxa de sucesso em encontrar uma solução válida;
    \item número de gerações até a solução;
    \item tempo total de execução;
    \item distribuição dos critérios de parada (ótimo encontrado, limite de gerações, limite de tempo);
    \item impacto do controle por tempo como critério adicional de parada.
\end{enumerate}

Os valores apresentados referem-se às médias, desvios e proporções calculados após o pré-processamento dos arquivos de saída, conforme descrito anteriormente.

\subsection{Taxa de Sucesso}

A Tabela~\ref{tab:sucesso} resume as taxas de sucesso das quatro estratégias para os diferentes tamanhos do tabuleiro considerados. Os resultados mostram que:

\begin{itemize}
    \item O \textbf{elitismo} apresentou, em geral, a maior taxa de sucesso, especialmente para tabuleiros maiores.
    \item A \textbf{substituição geracional} manteve um desempenho estável, porém inferior ao elitismo em cenários mais complexos.
    \item O esquema de \textbf{população fixa simples} teve desempenho equivalente ao geracional em tabuleiros menores, mas reduziu sua taxa de sucesso à medida que $n$ aumentou.
    \item O \textbf{controle por idade (GAVaPS)} demonstrou desempenho competitivo, superando a população fixa em vários tamanhos de tabuleiro e aproximando-se do elitismo nos cenários intermediários.
\end{itemize}

Esses resultados reforçam achados clássicos na literatura: estratégias que preservam diversidade populacional tendem a ser mais robustas em problemas combinatórios de difícil convergência.

\subsection{Número de Gerações até a Solução}

Ao analisar o número médio de gerações necessárias para encontrar uma solução, observam-se padrões complementares:

\begin{itemize}
    \item O \textbf{elitismo}, apesar da alta taxa de sucesso, frequentemente precisou de mais gerações para convergir, caracterizando um processo de refinamento contínuo.
    \item A \textbf{substituição geracional} convergiu mais rapidamente em tabuleiros pequenos, mas perdeu eficiência à medida que a complexidade aumentou.
    \item A \textbf{população fixa simples} apresentou alta variação, sugerindo maior sensibilidade a escolhas aleatórias durante a evolução.
    \item O \textbf{GAVaPS} destacou-se por manter um bom equilíbrio entre profundidade de busca e velocidade de convergência, especialmente quando o limite de tempo passou a ser o critério de parada mais frequente.
\end{itemize}

Essas observações sugerem que o controle por idade ajusta dinamicamente a diversidade da população, evitando convergência prematura sem exigir a preservação rígida de elite.

\subsection{Tempo Total de Execução}

A inclusão de um \textbf{critério de parada baseado em tempo} revelou diferenças significativas entre as estratégias:

\begin{itemize}
    \item O \textbf{controle por idade} foi o menos afetado pelo limite de tempo, indicando boa progressão mesmo em populações mais heterogêneas.
    \item O \textbf{elitismo}, embora eficiente em termos de taxa de sucesso, teve execuções mais longas devido à necessidade de avaliar repetidamente indivíduos elite.
    \item A \textbf{substituição geracional} foi a mais estável, com tempos relativamente previsíveis.
    \item A \textbf{população fixa simples} apresentou encerramentos prematuros com maior frequência, especialmente em instâncias grandes.
\end{itemize}

Esses resultados mostram que o controle por tempo não apenas evita execuções excessivamente longas, mas também funciona como indicador indireto da “velocidade evolutiva” de cada estratégia.

\subsection{Critérios de Parada}

As análises mostraram três comportamentos distintos:

\begin{enumerate}
    \item Em tabuleiros pequenos, a maior parte das execuções parou por \textbf{solução encontrada}, independentemente da estratégia.
    \item Em tabuleiros intermediários, o elitismo manteve taxas altas de sucesso, enquanto os demais passaram a atingir mais frequentemente o limite de gerações.
    \item Em tabuleiros grandes, o \textbf{limite de tempo} tornou-se dominante para todas as estratégias, mas o GAVaPS manteve a menor proporção de encerramentos por timeout.
\end{enumerate}

Isso reforça a utilidade do controle por tempo como complemento ao controle por gerações em experimentos de grande porte.

\subsection{Desempenho Comparativo Geral}

Ao consolidar as métricas, observam-se três conclusões recorrentes:

\begin{itemize}
    \item O \textbf{elitismo} continua sendo o método mais consistente em termos de taxa de sucesso, ao custo de maior tempo de execução.
    \item A \textbf{substituição geracional} e a \textbf{população fixa simples} são mais rápidas, porém menos robustas.
    \item O \textbf{controle por idade}, inspirado na abordagem GAVaPS, oferece o melhor compromisso entre diversidade, velocidade e capacidade de adaptação ao limite temporal.
\end{itemize}

Como resultado, embora nenhuma estratégia seja dominadora em todos os cenários, o controle por idade sobressai como alternativa promissora em problemas combinatórios intensivos, especialmente quando há interesse em limitar o tempo computacional.

\subsection{Considerações Sobre o Controle de Tempo}

A adição do tempo como critério de parada alterou o equilíbrio entre as estratégias. Em particular:

\begin{itemize}
    \item Sistemas com \textbf{população dinâmica} (como o controle por idade) respondem melhor ao limite temporal, pois a renovação constante evita estagnação.
    \item Estratégias que preservam explicitamente indivíduos muito aptos sofrem com ciclos mais longos de avaliação.
    \item O critério de tempo melhora a comparabilidade entre cenários, reduzindo assimetrias na duração dos experimentos.
\end{itemize}

Essa integração aumenta a utilidade prática dos AGs, aproximando o estudo de condições reais, onde recursos computacionais são limitados.

\subsection{Síntese dos Resultados}

Com base nas análises, os principais achados são:

\begin{itemize}
    \item O controle por idade demonstrou excelente equilíbrio entre desempenho e custo computacional.
    \item O elitismo apresentou as maiores taxas de convergência, porém com maior tempo médio por execução.
    \item A combinação de limite de gerações e limite de tempo se mostrou fundamental para garantir comparabilidade e previsibilidade.
    \item Os resultados confirmam que políticas de substituição têm impacto decisivo na eficiência evolutiva de AGs.
\end{itemize}

Esses aspectos fundamentam as conclusões apresentadas na próxima seção e sustentam o argumento de que diferentes esquemas de controle populacional podem ser mais adequados dependendo das restrições do problema e do ambiente computacional disponível.
